\begin{papercard}{zou2019}
\cf{Stages} \stagetag{C} \stagetag{I} \stagetag{H} (also \stagetag{A}). A content analysis of 161 consumer data-breach notifications covering readability metrics, structure, risk communication (hedge terms, ``no evidence'' claims) and presentation of protective actions. A three-coder codebook pilot reached Fleiss' $\kappa=.75$. It is a model for coding \emph{risk language} in privacy texts.

\cf{Aim} To examine readability and usability problems that may impair consumers' understanding of breach risks and available protective measures, a precondition for acting \pc{zou2019}{4}{We focus on examining the readability and usabil}.

\cf{Corpus and sampling} Source: the Maryland Attorney General's public database, chosen as the largest of the five public state databases and for its metadata (cause, data types) \pc{zou2019}{4}{Of the five states, Maryland’s database includes the largest number of breach notifications}. All 548 notices from January to June 2018 were downloaded. After removing duplicates and non-consumer letters, 326 remained, and 161 (about 50\,\%) were randomly selected, with the list in an appendix \pc{zou2019}{4}{In total, we obtained 548 data breach notifications}. Rules for multiple versions: analyse the version with the larger audience (adults over guardians, consumers over employees), and the first version by default \pc{zou2019}{4}{We analyzed the version with a potentially larger audience}. The sample was cross-referenced with the Privacy Rights Clearinghouse to gauge magnitude \pc{zou2019}{5}{We also cross-referenced our sample with the Privacy Rights Clearinghouse}.

\cf{Quantitative analysis} FRES, Flesch--Kincaid grade and Gunning Fog via readable.io, plus word and sentence counts \pc{zou2019}{5}{We also looked into the Gunning Fog index (FOG)}. Reading time was estimated at 250 words per minute, following McDonald and Cranor \pc{zou2019}{6}{we assumed a reading speed of 250 words per minute}.

\cf{Qualitative coding} One researcher built an initial codebook by thematic coding and affinity diagramming. Three researchers independently coded a random 20 notices (12.4\,\%), reconciled and revised, reaching Fleiss' $\kappa=.75$. The final codebook had 9 categories, 38 codes and 136 sub-codes (e.g.\ ``absolutely'', ``maybe'', ``no'', ``no evidence''). The corpus was then split and single-coded \pc{zou2019}{5}{Three members of the research team then independently}, \pc{zou2019}{5}{The final codebook (see Appendix B in auxiliary ma}.

\cf{Findings} \emph{Readability}: median FRES 46.7, with 97\,\% ``fairly difficult'' or ``difficult''; median grade 10.0 \pc{zou2019}{5}{This means that 97\% of the notifications were fairly difficult or difficult to read}. Median 1,575 words, about 6.3 minutes. \emph{Structure}: 67\,\% used headings, and 94\,\% of those followed California's template wording \pc{zou2019}{6}{100 (94\%) followed California’s wording and order requirements}. \emph{Risk communication}: hedge terms appeared in 70\,\% \pc{zou2019}{7}{Hedge terms such as “maybe” and “likely” were used in 112 (70\%) notifications}. ``No evidence of misuse'' claims appeared in 40\,\%, and lack of evidence is not evidence of absence \pc{zou2019}{7}{lack of evidence is not evidence of absence of harm}. Only 9 notices linked data types to concrete misuse scenarios, the effective strategy \pc{zou2019}{7}{connecting the types of breach infor}. Only 64\,\% stated when the breach occurred (``uninformed exposure time''). \emph{Actions}: median 8 recommended actions, risking choice overload, with strong actions such as a credit freeze hidden in appendices (73\,\%) \pc{zou2019}{8}{118 (73\%) notifications listed it as one of many options in an appendix}.

\cf{Limitations stated} One state's database, no direct evidence of effects on behaviour, and hypotheses that need user studies \pc{zou2019}{9}{our content analysis does not provide direct evidence of how text and format}.

\cf{Recommendations} Plain, concise language; prioritised actions with reasons; avoid hedge terms and ``no evidence'' defences or pair them with warnings; visual emphasis on deadlines; measurable readability standards; multiple delivery channels \pc{zou2019}{10}{we argue that companies should avoid “no evidence” claims and hedge terms}.

\cf{Use for this thesis} A directly transferable coding approach for \emph{how apps talk about risk}. Code certainty language (hedges such as ``may'' or ``might'' vs.\ definite statements), reassurance claims, and whether data types are linked to concrete consequences. The paper's pilot design (three coders, a random 12\,\% subset, Fleiss' $\kappa$ reported before splitting the corpus) is a clear precedent for the codebook pilot, and its limitation (single coding after the pilot) should be named if repeated. The readability metrics and the 250 wpm reading-time convention are standard, but German texts need German-specific formulas (e.g.\ Amstad's FRES adaptation), so record that as a decision.
\end{papercard}
